
\documentclass{article}
\usepackage{colortbl}
\usepackage{makecell}
\usepackage{multirow}
\usepackage{supertabular}

\begin{document}

\newcounter{utterance}

\centering \large Interaction Transcript for game `cladder', experiment `full\_v1.5\_default', episode 4345 with Qwen3.6{-}27B.
\vspace{24pt}

{ \footnotesize  \setcounter{utterance}{1}
\setlength{\tabcolsep}{0pt}
\begin{supertabular}{c@{$\;$}|p{.15\linewidth}@{}p{.15\linewidth}p{.15\linewidth}p{.15\linewidth}p{.15\linewidth}p{.15\linewidth}}
    \# & \multicolumn{2}{c}{Player} && \multicolumn{2}{c}{Game Master} \\
    \hline

    \theutterance \stepcounter{utterance}  
    & & & \multicolumn{4}{p{0.6\linewidth}}{
        \cellcolor[rgb]{0.9,0.9,0.9}{
            \makecell[{{p{\linewidth}}}]{
                \texttt{\tiny{[P1$\langle$GM]}}
                \texttt{You are an expert in causal inference. The following question is not a typical commonsense query, but rather a meticulously designed question created by a professor specializing in causal inference, intended to assess the students' mastery of the course content.} \\
\texttt{Imagine a self{-}contained, hypothetical world with only the following conditions, and without any unmentioned factors or causal relationships: Maternal smoking status has a direct effect on infant's birth weight and freckles. Health condition has a direct effect on infant's birth weight and freckles. Infant's birth weight has a direct effect on freckles. Health condition is unobserved. Method 1: We look directly at how maternal smoking status correlates with freckles in general. Method 2: We look at this correlation case by case according to infant's birth weight. To understand how maternal smoking status affects freckles, is it more correct to use the Method 1 than Method 2?} \\
\texttt{Start your answer with 'Yes' or 'No', followed by additional reasoning or evidence to support your explanation.} \\
            }
        }
    }
     \\ \\

    \theutterance \stepcounter{utterance}  
    & \multicolumn{4}{p{0.6\linewidth}}{
        \cellcolor[rgb]{0.561,0.737,0.561}{
            \makecell[{{p{\linewidth}}}]{
                \texttt{\tiny{[P1$\rangle$GM]}}
                \texttt{Yes} \\
\\ 
\texttt{**Reasoning:**} \\
\\ 
\texttt{To determine the correct method, we must analyze the causal structure (Directed Acyclic Graph, or DAG) implied by the problem statement and identify whether conditioning on specific variables introduces bias.} \\
\\ 
\texttt{**1. Constructing the Causal Graph:**} \\
\texttt{Based on the description, the variables and their direct causal relationships are:} \\
\texttt{*   \$S\$ (Maternal Smoking) \$\textbackslash{}rightarrow\$ \$W\$ (Birth Weight)} \\
\texttt{*   \$S\$ (Maternal Smoking) \$\textbackslash{}rightarrow\$ \$F\$ (Freckles)} \\
\texttt{*   \$H\$ (Health Condition) \$\textbackslash{}rightarrow\$ \$W\$ (Birth Weight)} \\
\texttt{*   \$H\$ (Health Condition) \$\textbackslash{}rightarrow\$ \$F\$ (Freckles)} \\
\texttt{*   \$W\$ (Birth Weight) \$\textbackslash{}rightarrow\$ \$F\$ (Freckles)} \\
\\ 
\texttt{Note that \$H\$ is an unobserved confounder for the relationship between \$S\$ and \$F\$ (since \$H\$ affects both, though the prompt says \$S\$ has a *direct* effect on \$F\$, \$H\$ is a common cause of \$F\$ and \$W\$, and \$S\$ causes \$W\$. Wait, let's look closer at confounding for the \$S \textbackslash{}rightarrow F\$ effect).} \\
\\ 
\texttt{Actually, let's look at the paths between \$S\$ (Smoking) and \$F\$ (Freckles):} \\
\texttt{1.  **Direct Path:** \$} \\
            }
        }
    }
    & & \\ \\

    \theutterance \stepcounter{utterance}  
    & & & \multicolumn{2}{p{0.3\linewidth}}{
        \cellcolor[rgb]{0.9,0.9,0.9}{
            \makecell[{{p{\linewidth}}}]{
                \texttt{\tiny{[GM$|$GM]}}
                \texttt{Yes} \\
            }
        }
    }
    & & \\ \\

    \theutterance \stepcounter{utterance}  
    & & & \multicolumn{2}{p{0.3\linewidth}}{
        \cellcolor[rgb]{0.9,0.9,0.9}{
            \makecell[{{p{\linewidth}}}]{
                \texttt{\tiny{[GM$|$GM]}}
                \texttt{yes} \\
            }
        }
    }
    & & \\ \\

    \theutterance \stepcounter{utterance}  
    & & & \multicolumn{2}{p{0.3\linewidth}}{
        \cellcolor[rgb]{0.9,0.9,0.9}{
            \makecell[{{p{\linewidth}}}]{
                \texttt{\tiny{[GM$|$GM]}}
                \texttt{game\_result = WIN} \\
            }
        }
    }
    & & \\ \\

\end{supertabular}
}

\end{document}
